\section{Introducción}
\label{sec:intro}

Los avances recientes en multimodal large language models (MLLMs) han ampliado drásticamente el alcance del razonamiento visual. El vision--language instruction tuning ahora permite que un único backbone responda preguntas abiertas sobre secuencias de video largas \cite{llava,openai2023gpt4v,qwen2.5}, mientras que técnicas de extensión de contexto como FlashAttention-2 y RingAttention llevan la ventana efectiva al régimen del millón de tokens \cite{dao2023flashattention2,ringattention,peng2024yarn}. Estos avances sustentan una nueva generación de {\em asistentes de video en streaming} y \emph{robots humanoides} que prometen comprensión continua y en tiempo real de la escena en teléfonos móviles, gafas AR y robots edge~\cite{deepmind2023astra,openai2023gpt4o,waisberg2024meta,MorganStanley2024FigureAI}.

Streaming video understanding (SVU) difiere de la comprensión convencional de video offline (OVU). Los modelos offline ven el clip completo y la consulta del usuario antes de la inferencia, por lo que pueden adaptar cada paso de compresión o recuperación. En streaming, los frames llegan de manera incremental y las consultas futuras son desconocidas, lo que obliga a que todo el procesamiento previo a la consulta sea {\em query-agnostic}. Además, la memoria del dispositivo es fija, pero el Transformer emite cientos de tokens por frame, por lo que la caché key--value (KV) crece linealmente. Por ejemplo, un clip de 15 minutos a 10 fps procesado por LLaVA-Next-Video-7B ya requiere \(\sim\!100\)\,GB de almacenamiento KV, muy por encima de las decenas de gigabytes disponibles en plataformas móviles o robóticas \cite{zhang2024llavanext-video,Karumbunathan2022JetsonAGXOrin}.

Los trabajos previos abordan las restricciones de memoria en videos largos en tres etapas (Fig.\ref{fig:background}). {\em Frame Sampling}~\cite{feichtenhofer2019slowfastnetworksvideorecognition} descarta frames antes de la codificación, reduciendo memoria pero degradando severamente la cobertura temporal y la precisión. {\em Input-Vision Compression (IVC)}~\cite{shen2024longvu,tao2024dycokedynamiccompressiontokens} elimina vision tokens redundantes después de la codificación, reduciendo la carga de Prefill pero requiriendo todavía que el conjunto completo de vision tokens se almacene en memoria. {\em KV cache Compression (KVC)}~\cite{li2024snapkv,fu2024headkv} selecciona tokens relevantes para la consulta después del paso de Prefill, ofreciendo la mayor precisión, pero solo después de materializar la caché KV completa. El desafío se intensifica en escenarios de streaming: el uso de memoria de Frame Sampling, IVC y KVC crece casi linealmente con la longitud del video, hasta superar los límites del dispositivo. El KV cache offloading (por ejemplo, ReKV~\cite{Shangzhe2025streaming}) amplía el espacio de memoria, pero incurre en costosas transferencias de datos repetidas para cada consulta. Por lo tanto, ningún enfoque existente ofrece la propiedad clave que SVU necesita: una compresión de video en streaming \textit{independiente de la longitud} y \textit{query-agnostic}.


Un enfoque natural para abordar las restricciones de memoria en streaming video es explotar la fuerte \textit{redundancia espaciotemporal} de los flujos de video. Introducimos \textbf{InfiniPot-V}, el primer framework diseñado específicamente para SVU con memoria restringida. InfiniPot-V es \textit{training-free}, \textit{query-agnostic} y opera de forma \textit{continua} durante la inferencia. Cuando la caché KV alcanza un umbral de memoria definido por el usuario \(M\), realiza una compresión in-place que libera espacio para nuevos frames preservando la esencia semántica del contexto previo. Esta compresión está guiada por dos métricas livianas y complementarias. \textit{Temporal-axis Redundancy (TaR)} modela los Key embeddings como un tensor 3D en el tiempo y elimina tokens con alta similitud coseno con frames recientes, podando de forma efectiva contenido estático o repetitivo. \textit{Value-Norm Importance (VaN)} ordena los tokens restantes según la norma \(\ell_2\) de sus vectores Value ---un fuerte proxy agnóstico al modelo de saliencia semántica--- y aplica una estrategia de pooling adaptativa por capa. Esta compresión es altamente eficiente, añade una latencia despreciable y a la vez hace cumplir estrictamente los límites de memoria.

Una evaluación extensa confirma la efectividad de este diseño. En cuatro MLLMs open-source (3B y 7B) y seis benchmarks de video largo ---que cubren tanto tareas offline (VideoMME, EgoSchema, MLVU, LongVideoBench) como de streaming (RVS-Ego/Movie, OVO-Bench, StreamingBench)--- InfiniPot-V reduce el uso de longitud de contexto de entrada hasta 6K para contextos de 50K tokens, con una precisión que iguala o supera a las líneas base de caché completa. Mantiene rendimiento en tiempo real a 14 frames por segundo con solo 0.5\% de overhead de compresión. Además, su naturaleza query-agnostic ofrece beneficios claros en configuraciones de diálogo multi-turno (Appendix.~\ref{appn:multi_turn}). Al eliminar el cuello de botella de la caché KV sin reentrenamiento ni dependencia de la consulta, InfiniPot-V allana el camino para asistentes multimodales prácticos ejecutados en dispositivo.
